% ============================================================================
\section{Discussion and conclusion}\label{sec:conclusion}
% ============================================================================

\subsection{Managerial implications}

The study asked how contextual forecasts should be turned into replenishment decisions when perishable stock ages and deliveries share vehicles. Five implications follow for planners of fresh-food distribution.

\begin{enumerate}[leftmargin=*]
\item \textbf{Start from a stock-age-based replenishment rule, tune it finely, and benchmark against it before integrating.} Daily replenishment at a high service quantile, with the orders routed optimally, was cheaper than every integrated inventory-routing policy in all settings but the one with expensive routing, by 2.0--11.9\%, across three network layouts, shelf lives of two to four days, lower and higher lost-sales penalties, tighter and looser fleets, more variable demand, partial LIFO issuing and observed demand of a fresh-food retailer; the margin was smallest when lost sales were cheap, so that high service is worth less. In these settings the travel saved by consolidating deliveries is smaller than the extra holding, waste and lost sales that consolidation causes. A finer grid for the service level and the visit trigger alone lowered the cost of the deliver-up-to-level rule by up to 0.6\%. This extends the finding of \citet{haijema2019improved} and \citet{crama2018stochastic} that well-tuned simple rules are hard to beat to a setting with shared routes and contextual forecasts.
\item \textbf{Learn the order-up-to level, not only the forecast, but check that what is learned carries over.} Training the order-up-to level of the deliver-up-to-level rule on the realised cost of its decisions reduced cost by about 1\% in nine of thirteen settings and on fresh test episodes of the base case, where it reached almost the same fill rate with 28\% less waste, which no re-tuning of the rule's service level and trigger achieves. The gain was smallest, and not significant, where the rule already wasted little (four-day shelf life), with a loose fleet, and where demand was too variable for the context to anticipate it. On observed demand of new products the learned corrections did not carry over, and the rule was cheaper: corrections should be validated on products and periods like those on which they will be used, and the rule kept where they are not. The learning also needs the rule: the same network trained to order directly was more expensive than UL-DFL in twelve settings. The policy is operated exactly like the rule: it remains a rule with a visit trigger and an order-up-to level that planners can inspect. The learned corrections translate into practice that does not require machine learning: shift safety stock from the days after the weekly peak to the days before it, and, where promotions widen the forecast spread, lower the order-up-to level on promotion days.
\item \textbf{Integration pays when routing is expensive, and then the safety stock of the plan matters most.} With three times higher travel cost, a lookahead with day-specific quantile levels was the cheapest policy, level with UL-DFL, and every uncertainty-aware lookahead beat daily base-stock. Scenario-based lookahead with the scenario set of \citet{jia2025scenario}, which does not extend beyond the 0.9 quantile of the forecasts, under-stocked; with scenarios drawn from bands of the full quantile grid, which reach further into the upper tail, it would have been the cheapest policy in that setting. Deterministic lookahead with mean forecasts was the worst policy in every setting.
\item \textbf{If a deterministic plan is used, its safety stock should depend on the planning day.} Decision-focused learning on top of the lookahead learned to keep a high service quantile only for the day that is actually executed and to plan later days near the 0.8 quantile. A lookahead with these two quantile levels, tuned on validation data, removes about half of the waste of the single-quantile plan in the base case and was at least as cheap as the learned targets in all thirteen settings.
\item \textbf{Look at store-level service, not only at averages.} Integrated deterministic plans concentrated stock-outs at the store whose visit costs the most travel per unit sold, whereas the inventory rules served all stores alike. Retailers with per-store service targets should constrain or monitor this effect.
\end{enumerate}

\subsection{Limitations and further research}

The synthetic demand is generated by a documented contextual model; the real-data case uses observed sales of one retailer, from which latent demand had to be reconstructed under stock-outs and rescaled to the network's capacities, and from which products whose level shifted strongly were excluded. The test episodes of the synthetic configurations had been evaluated in an earlier round of the study; the primary hypotheses are therefore tested on fresh test episodes of the base case, and the other synthetic configurations are secondary evidence. In the real-data case, model selection on validation episodes did not protect against corrections that failed to transfer to new products; regularising the corrections towards the rule, or validating across products and periods, are natural remedies. Exact routing by enumeration limits the network to about a dozen stores; larger networks require heuristic routing, which would weaken the clean separation between replenishment and routing quality used here. Lead time is zero and each store sells one product; positive lead times would increase the value of stock-age information \citep{broekmeulen2009heuristic}, and multiple products would couple stores through shared vehicle space. The lookahead policies relax routing integrality beyond the implemented day, which favours visiting stores more often than a fully integer lookahead would. ScenPO-SS was implemented with the scenario set of the original study; its cost depends on how far the scenarios reach into the upper tail of the forecast, and tuning this choice on validation data, as done for the quantile levels of the deterministic lookahead, could make it more competitive. UL-DFL learns the order-up-to level but keeps the rule's visit trigger and sees fleet capacity only as a proportional scaling during training; learning the trigger as well, or training against the exact routing cost, could extend its advantage to settings with expensive routing. LA-DFL was trained on the decision states visited by LA-Quantile in one round; iterating between training and data collection, or training the scenario generator of ScenPO-SS decision-focusedly, are natural extensions. Finally, deep reinforcement learning has shown promise for single-location perishable inventory \citep{demoor2022reward,temizoz2025deep}; extending it to the routing-coupled action space studied here is an open problem.

\subsection{Conclusion}

Replenishing perishables on shared routes is a perishable inventory problem first and a routing problem second. Under a common closed-loop simulator with exact optimisation, stock-age-based replenishment rules were the strongest policies without learning unless travel was expensive, in which case a lookahead with day-specific safety quantiles matched them. Decision-focused learning made a well-tuned rule about 1\% cheaper and markedly less wasteful in most settings, including fresh test episodes, by learning from realised cost when to deviate from a fixed service quantile; on observed demand of new products its corrections did not carry over. On top of a deterministic lookahead it revealed why such plans mis-stock, a lesson that a day-specific quantile level implements without learning. The structural results, the experimental protocol and the code are available for testing further policies on the same benchmark.
